Această lucrare prezintă proiectarea și implementarea aplicației \project, un sistem web care analizează anunțuri de vânzare pentru autoturisme second-hand și estimează dacă prețul cerut este corect, prea mic sau prea mare. Aplicația pornește de la linkuri \autovit\ sau \mobilede, extrage informații relevante despre vehicul, normalizează datele, aplică modele de regresie antrenate pe date colectate automat și întoarce un verdict explicabil utilizatorului.

Lucrarea descrie întregul flux tehnic: colectarea datelor prin web scraping, dificultățile generate de site-urile pe care s-a realizat scraping și mecanisme anti-automatizare, deduplicarea anunțurilor, transformarea datelor brute JSON în seturi CSV pregătite pentru învățare automată, selecția caracteristicilor folosind teste statistice și antrenarea mai multor modele de regresie. Au fost comparate modele SVR, Ridge, KNN, Random Forest, Extra Trees, Gradient Boosting și un ansamblu de votare, antrenate pe ținta logaritmică a prețului. Pe setul combinat extins folosit pentru evaluare, cel mai bun model a obținut un scor \(R^2\) de aproximativ \(0.923\), iar erorile absolute medii au fost în intervalul \(2696\)-\(3320\) EUR, în funcție de model.

Rezultatul practic este o aplicație FastAPI cu interfață web în limba română, autentificare, istoric pe cont, afișarea anunțului, estimări pe model, medie ponderată configurabilă și recomandări de anunțuri similare. Lucrarea evidențiază atât utilitatea unei astfel de soluții pentru cumpărători, cât și limitele principale ale abordării cu Web Scraping si Machine Learning: dependența de calitatea datelor colectate, instabilitatea structurii site-urilor, acoperirea redusă pentru anumite categorii de mașini și riscul ca anunțurile foarte atipice să fie estimate mai slab.
